\title{Parameter-Efficient Orthogonal Finetuning via Factorization}

\begin{abstract}
The prevalence of large foundation models has grown significantly, yet their immense computational cost makes training them from the ground up largely impractical. This creates a strong incentive to find effective strategies for adapting these models to new tasks. In this work, we investigate Orthogonal Finetuning (OFT), a theoretically motivated approach for adapting models to downstream tasks. Although OFT tends to generalize well, it still requires a substantial number of trainable parameters, as orthogonal matrices are intrinsically high-dimensional. To enhance parameter-efficiency, we analyze OFT from the standpoint of information transfer and articulate several guiding principles for reducing parameter count. Drawing inspiration from the efficiency of the Cooley-Tukey fast Fourier transform algorithm in transmitting information, we introduce a novel, efficient orthogonal parameterization based on butterfly matrices. Integrating this structure into OFT, we develop a new finetuning methodology termed Orthogonal Butterfly (BOFT). BOFT generalizes OFT by encompassing it as a special case, thereby presenting a unified framework for orthogonal finetuning. Our extensive experiments adapt large-scale vision transformers, language models, and text-to-image diffusion models for a variety of downstream visual and linguistic tasks, demonstrating the efficacy of our approach.
\end{abstract}

\section{Introduction}

Recent advances with models such as ChatGPT~\cite{brown2020language,chen2021evaluating} and Stable Diffusion~\cite{rombach2022high} have highlighted the exceptional generalization capabilities of large foundation models. This leap in performance has come at the cost of exponentially more parameters 

(\emph{e.g.}, GPT-3 contains approximately 175B parameters), rendering the process of training such models from the ground up increasingly inaccessible. Consequently, meaningful advances increasingly depend on the capacity to modify existing foundation models rather than rebuilding them. Efficient adaptation of these models to downstream applications is, therefore, essential. Typically, three main strategies are used for adapting models efficiently: \emph{model finetuning}~\cite{hulora2022,zaken2022bitfit,guo2020parameter,raffel2020exploring,qiu2023controlling,zhang2022adaptive,chavan2023one}, which tunes a selected portion of parameters without altering the original architecture; \emph{adapter tuning}~\cite{rebuffi2017learning,houlsby2019parameter,pfeiffer2020adapterfusion,lin2020exploring,he2021towards}, where trainable modules are integrated into the existing network; and \emph{prompt tuning}~\cite{li2021prefix,lester2021power}, which attaches tunable token embeddings to the input sequence. Among these approaches, model finetuning is particularly appealing for its simplicity, strong performance, and the fact that it does not incur additional inference delay.

The central concept underlying model finetuning is to maintain a degree of similarity between the pretrained and finetuned models, as measured by certain metrics, so that the knowledge gained during pretraining is not lost. Traditional finetuning approaches often rely on using a small learning rate, which serves as an informal mechanism to limit divergence between the original and finetuned model. Given the inherently structured composition of weight matrices, a more systematic strategy involves safeguarding the relational properties within these matrices, specifically preserving the pairwise angular relationships between neurons. Building on this intuition, the Orthogonal Finetuning~(OFT)~\cite{qiu2023controlling} paradigm has emerged. In OFT, all neurons within a particular layer are transformed via a shared orthogonal matrix, thereby guaranteeing that the angles between neurons remain unchanged throughout the finetuning phase. While OFT has achieved notable results in terms of both generalization and convergence when applied to the finetuning of text-to-image diffusion models~\cite{qiu2023controlling}, the high dimensionality of orthogonal matrices leads to a considerable increase in trainable parameters. To mitigate this issue, OFT leverages a block-diagonal formulation that reduces parameter count. Nevertheless, this approach introduces certain drawbacks: the orthogonal matrices acquire a rigid sparsity structure and are applied separately to each block, resulting in parameter efficiency at the cost of expressiveness. The block-diagonal design may impose restrictive inductive biases (\emph{e.g.}, diminished expressiveness and an inability to emulate traditional linear transformations), and crucially, the optimal way to determine a sparsity pattern for such matrices is still unresolved.

Addressing this issue hinges on creating a dense orthogonal matrix in a parameter-efficient manner. Although the conventional approach would require $\mathcal{O}(d^2)$ parameters for a $d$-dimensional dense orthogonal matrix, we introduce an alternative strategy: constructing such a matrix by combining several factorized sparse matrices. This method capitalizes on the observation that computational cost can be exchanged for a reduction in parameter count. Because the orthogonal matrix is expressed as the product of sparse matrices, the multiplication process must be executed multiple times during each training step. Framing matrix factorization from a new angle, we view the construction of a dense orthogonal matrix as a problem of information propagation. Within this perspective, assembling a dense orthogonal matrix through sparse matrix products is analogous to conveying information across a grid-like network. This viewpoint affords the flexibility to design a range of sparse matrix factorizations, all of which minimize the number of learnable parameters without compromising the ability to represent dense matrices.

For enhanced parameter efficiency under this information propagation framework, we draw on the butterfly architecture of the Cooley-Tukey fast Fourier transform algorithm~\cite{cooley1965algorithm}, which demonstrates efficient data transfer among various nodes~\cite{klasing1994broadcasting}. In particular, the butterfly graph underpinning the Cooley-Tukey algorithm suggests a specialized method for sparse matrix decomposition—referred to as butterfly factorization. Through this approach, a $d$-dimensional dense matrix is constructed from the product of $\mathcal{O}(\log d)$ sparse matrices, each containing only $\mathcal{O}(d)$ nonzero elements. When these sparse components are constrained to be orthogonal, we achieve an efficient orthogonal matrix representation with just $\mathcal{O}(d\log d)$ parameters—substantially fewer than required in the naive case. Leveraging this compact orthogonal parameterization, we introduce a new, highly parameter-efficient finetuning technique called Orthogonal Butterfly~(BOFT). Notably, BOFT generalizes OFT as a special case and offers a unified framework for orthogonal finetuning. Both the block-diagonal and butterfly structures share the trait of sparsity, which imparts parameter efficiency by incorporating structured zeros into orthogonal matrices. Comparing our approach with the low-rank matrix construction in LoRA~\cite{hulora2022}, we note that both low-rank and sparse matrices form fundamental classes of structured matrices~\cite{candes2011robust} that realize parameter savings.

In contrast to the block-diagonal paradigm employed by OFT to balance the trade-off between model expressiveness and regularity, BOFT leverages the butterfly structure to achieve a more continuous transition between matrices drawn from the full orthogonal group (for expressivity) and identity matrices (for regularization purposes). This structural flexibility allows us to effectively constrain the hypothesis space to a narrower, more relevant subset of the orthogonal group, enhancing performance on downstream tasks. Given the prevalence of the butterfly pattern in a wide array of efficient linear algorithms—including the discrete Fourier and discrete cosine transforms—we contend that our structured parameterization inherently induces an inductive bias conducive to better generalization and reduced risk of overfitting. This perspective is motivated by the fact that the butterfly architecture is capable of exactly reproducing numerous well-established linear transforms, a feat unattainable by the block-diagonal design utilized in OFT. Our primary contributions are summarized as follows:

\begin{itemize}[leftmargin=*,nosep]
    \item We introduce a new perspective on parameter efficiency for orthogonal finetuning by framing the problem within an information transmission context. Specifically, we recast the task of learning a compact yet dense orthogonal matrix as an information transmission challenge on a graph with grid topology.
    \item Drawing inspiration from the butterfly structure underlying the Cooley-Tukey algorithm, we develop Orthogonal Butterfly, a parameter-efficient approach for orthogonal finetuning situated within our information transmission paradigm.
    \item We offer several theoretical justifications for why BOFT achieves substantial reductions in trainable parameters while preserving high levels of expressivity and training robustness. Moreover, the matrix decomposition intrinsic to BOFT provides an appealing mechanism for weight interpolation (see Figure~\ref{fig:boft_interpolation}).
    \item We pioneer the application of orthogonal finetuning across a broad spectrum of tasks that extend beyond controllable text-to-image generation~\cite{qiu2023controlling}, demonstrating BOFT’s versatility as a general finetuning tool. Our experiments span adaptation scenarios in both computer vision and natural language processing, with BOFT consistently outperforming current leading approaches, underscoring its advantages in parameter efficiency and generalization capability.
\end{itemize}

\textbf{Parameter-efficient finetuning (PEFT).} As foundation models (\emph{e.g.}, \cite{brown2020language,radford2021learning,rombach2022high,kirillov2023segment}) continue to expand in scale and capability, optimizing these models for specific downstream applications in a parameter-efficient manner has attracted significant attention~\cite{treviso2023efficient,ding2023parameter,lialin2023scaling}. Within the broad family of PEFT strategies~\cite{houlsby2019parameter,aghajanyan2020intrinsic,hulora2022,edalati2022krona,wang2022adamix,gheini2021cross,zaken2022bitfit,guo2020parameter,sung2021training,ansell2022composable,lester2021power,li2021prefix,vu2022spot,he2021towards,mao2021unipelt,karimi2021compacter,liu2022few,sung2022lst,chen2022parameter,jia2022visual,chen2022adaptformer,zhang2022neural,jie2023fact,lian2022scaling,luo2023towards}, approaches grounded in reparameterization~\cite{aghajanyan2020intrinsic,hulora2022,edalati2022krona} are most relevant for our study. In particular, LoRA~\cite{hulora2022} augments the pretrained weight matrix by incorporating the product of two low-rank matrices, achieving notable success on natural language tasks. Since LoRA imposes a uniform, fixed rank for all introduced matrices, various follow-up works~\cite{zhang2022adaptive,valipour2022dylora,zhang2023increlora} have devised methods to adjust the rank across layers dynamically, thus distributing the parameter budget more judiciously. The straightforwardness of low-rank matrix reparameterization has led to its widespread adoption~\cite{dettmers2023qlora,zi2023delta,chavan2023one}. Drawing inspiration from the use of hyperspherical energy to understand generalization~\cite{liu2018learning,liu2021orthogonal}, \cite{qiu2023controlling} introduce orthogonal finetuning as an alternative mechanism to refine text-to-image diffusion models, wherein an orthogonal matrix is trained to transform neurons within the same layer. This approach consistently yields improved generalization and more stable convergence compared to LoRA. However, OFT generally entails a higher number of tunable parameters relative to LoRA, prompting the need to enhance its parameter efficiency. Furthermore, it remains to be investigated whether OFT can extend its utility to broader adaptation scenarios beyond text-to-image diffusion models~\cite{qiu2023controlling}. BOFT enhances the parameter efficiency of OFT by leveraging butterfly factorization, thereby allowing us to validate the effectiveness of orthogonal finetuning across diverse adaptation tasks.

\textbf{Butterfly structures.} The radix-2 Cooley-Tukey algorithm systematically breaks down the $N$-point discrete Fourier transform into a pair of $\frac{N}{2}$-point transforms in a recursive manner, yielding a distinctive butterfly structure which can be formally expressed as a product of several sparse matrices—collectively referred to as a butterfly matrix. These butterfly matrices have been previously exploited to parametrize orthogonal matrices, thereby sidestepping the need for pivoting in Gaussian elimination and improving computational efficiency~\cite{parker1995random}, in addition to stabilizing recurrent neural network training~\cite{jing2017tunable}, and enhancing kernel approximations~\cite{munkhoeva2018quadrature}. Other works such as~\cite{dao2019learning,dao2019kaleidoscope} leverage butterfly parameterizations to devise efficient linear transformations, while~\cite{chen2022pixelated,dao2022monarch} employ butterfly matrices to facilitate scalable training of neural networks. Moreover, butterfly structures have applications in rapid matrix-vector multiplication~\cite{michielssen1996multilevel,o2010algorithm}, data-sparse matrix approximations~\cite{li2015butterfly}, and transmission in networked systems~\cite{klasing1994broadcasting,li2003linear,rajasingh2016transmission}. In our work, in contrast to previous uses, we utilize butterfly structures specifically to improve the parameter efficiency of OFT.

\section{An Information Transmission View on Orthogonal Finetuning}

We begin by reviewing some foundational concepts in OFT. OFT performs finetuning of a pretrained weight matrix by expressing the updated matrix as a product of a learnable orthogonal matrix and the fixed pretrained weights. Unlike LoRA, which applies a low-rank additive perturbation to the original weights, OFT instead multiplies the pretrained matrix by an orthogonal matrix for updating. LoRA’s parameter efficiency arises from its low-rank design; in contrast, classic OFT attains efficiency through a block-diagonal orthogonal structure~\cite{qiu2023controlling}, where decreasing the block size enhances parameter efficiency. Figure~\ref{fig:comp_lora} provides a visual comparison between these strategies. The primary rationale behind leveraging an orthogonal transformation for weight finetuning lies in preserving inter-neuronal angular relationships~\cite{liu2018learning,liu2021orthogonal,qiu2023controlling}, thereby maintaining much of the semantic structure learned during pretraining. Specifically, OFT learns an orthogonal matrix $\bm{R}\in\mathbb{R}^{d\times d}$ corresponding to a pretrained linear transformation $\bm{W}^0\in\mathbb{R}^{d\times n}$, altering the original forward computation $\bm{z}=(\bm{W}^0)^\top\bm{x}$ to $\bm{z}=(\bm{R}\bm{W}^0)^\top\bm{x}$, where $\bm{x}\in\mathbb{R}^{d}$ and $\bm{z}\in\mathbb{R}^{n}$ represent the input and output vectors, respectively. For zero initialization, $\bm{R}$ is set to the identity matrix at the outset. To preserve the orthogonality of $\bm{R}$ during training, the Cayley parameterization, as proposed in~\cite{qiu2023controlling,liu2021orthogonal}, is adopted: $\bm{R}=(\bm{I}+\bm{Q})(\bm{I}-\bm{Q})^{-1}$, where $\bm{Q}$ is skew-symmetric ($\bm{Q}=-\bm{Q}^\top$). For enhanced parameter efficiency, OFT’s orthogonal matrix is structured in block-diagonal form, such that $\bm{R}=\text{diag}(\bm{R}_1,\bm{R}_2,\cdots,\bm{R}_r)$ with each block $\bm{R}_i\in\mathcal{R}^{b\times b}$ and $br=d$. The parameter savings afforded by this block-diagonal approach, however, are not without compromise; the method partitions the neuron's dimensions (i.e., columns of $\bm{W}^0$) into $r$ separate groups, with each segment independently transformed by a distinct orthogonal matrix. Although empirically effective, this partitioning based solely on index does not align with any intrinsic property of neuron representations, making dense orthogonal transformations potentially more suitable. This leads to a natural inquiry: \emph{Is it possible to devise a dense orthogonal matrix that retains the parameter-efficiency offered by the block-diagonal design?}

To tackle this issue, we suggest constructing a dense orthogonal matrix by multiplying several sparse orthogonal matrices. To create a coherent and easily graspable framework for analyzing the sparsity patterns in orthogonal matrix factorizations, we reinterpret the task of forming a dense orthogonal matrix within OFT as analogous to an information transmission problem. More precisely, forming a dense matrix $\bm{R}\in\mathbb{R}^{d\times d}$ through the product of $m$ square matrices, i.e., $\thickmuskip=2mu \medmuskip=2mu \bm{R}=\bm{B}_m\bm{B}_{m-1}\cdots\bm{B}_1$, can be visualized as information propagating through a $d\times (m+1)$ node grid, as depicted in Figure~\ref{fig:info_trans}. This interpretation is motivated by recognizing that a dense $d$-dimensional square matrix encodes comprehensive connections between two groups of $d$ nodes. For the matrix $\bm{R}$, a zero in the $(i,j)$ entry ($R_{ij}=0$) means that the $j$-th node’s information does not reach the $i$-th node, while a non-zero entry allows for this flow. Hence, decomposing the dense matrix $\bm{R}$ into a product $\bm{B}_m\bm{B}_{m-1}\cdots\bm{B}_1$ can be regarded as executing stepwise information transfers over the sequence of graphs each $\bm{B}_i$ represents. Information thus moves first via $\bm{B}_1$, progressing through to $\bm{B}_n$. Taking Figure~\ref{fig:info_trans} as an example, we analyze the decomposition $\bm{R}=\bm{B}_5\bm{B}_4\bm{B}_3\bm{B}_2\bm{B}_1$; its sparsity structures and the resulting graph are illustrated. The network shown in Figure~\ref{fig:info_trans} arises from expanding out the sequence of matrix multiplications, where each $\bm{B}_i$ is interpreted as encoding links from the $i$-th layer of nodes to the $(i+1)$-th layer. Specifically, entry $(j_1, j_2)$ of $\bm{B}_i$ represents the existence of a directed connection from node $j_2$ on the $i$-th layer to node $j_1$ on the $(i+1)$-th layer (with a value of zero indicating no such connection). For the product $\bm{B}_5\bm{B}_4\bm{B}_3\bm{B}_2\bm{B}_1$ to yield a dense matrix, every initial node must be able to reach each terminal node at the $6$-th layer. Restricting attention to $\bm{R}=\bm{B}_4\bm{B}_3\bm{B}_2\bm{B}_1$, mapping sources at the first layer to destinations at the $5$-th, we observe, for example, that information starting from node $1$ cannot be relayed to node $3$. This reflects that the compound sparsity of $\bm{B}_4\bm{B}_3\bm{B}_2\bm{B}_1$ includes a zero at the $(3,1)$ location. The same pattern arises for factorizations $\bm{R}=\bm{B}_3\bm{B}_2\bm{B}_1$ and $\bm{R}=\bm{B}_2\bm{B}_1$: the graph structure directly matches the sparsity layout. In general terms, for any dense $\bm{R}\in\mathbb{R}^{d\times d}$, the sequence of factors $\bm{B}_m,\cdots,\bm{B}_1$ should form a series of directed edges on a $d\times (m+1)$ lattice, with each edge connecting adjacent columns (levels) only, ensuring that each initial node can access all $(m+1)$-th level nodes via the transmission process.

Interpreting OFT through the lens of information transmission offers valuable insight into the sparsity characteristics of its factorization matrices. As depicted in Figure~\ref{fig:blk_diag}, the matrix $\bm{R}$ in the canonical OFT framework displays a block-diagonal arrangement. While this structure diminishes the count of trainable parameters, it falls short of achieving a fully dense configuration for $\bm{R}$. The objective is to assemble a dense orthogonal matrix by leveraging $m$ sparse orthogonal factors, striving to utilize the minimum number of essential trainable parameters.

From the information transmission standpoint, there are two principal requirements for achieving this goal: \emph{(i)} \textbf{dense connectivity}, ensuring every node at the initial layer can reach any node at the terminal layer via at least one connecting path, and \emph{(ii)} \textbf{minimum free edges}, where the cumulative edge count remains as low as feasible under the orthogonality constraint. Importantly, enforcing orthogonality imposes nuanced limitations on connections between adjacent layers. For instance, each factor $\bm{B}_i$ must be full-rank (an orthogonality prerequisite), necessitating $d$ edges to establish a bijective correspondence between nodes of the $i$-th and $(i=1)$-th layers, setting a lower bound of $d$ edges (\emph{e.g.}, four in the context of Figure~\ref{fig:info_trans}). These $d$ structural edges are vital for orthogonality, are untrainable (in a $d\times d$ orthogonal matrix with $d$ non-zero entries, these values are constrained to $\pm 1$), and therefore should not be included in the edge count.

Since orthogonal matrices inherently require fewer independent parameters than general dense matrices, the orthogonality constraint introduces intricate dependencies between possible edges. To illustrate, consider a $2\times 2$ orthogonal matrix: enforcing a zero in the $(1,1)$ entry compels a corresponding zero in the $(2,2)$ position—so one absent edge can dictate the absence of another. Accordingly, orthogonality may modify the number of admissible edges for certain connection patterns. Utilizing the information transmission viewpoint—especially by visualizing the sparsity mask of the constructed orthogonal matrix—proves advantageous for OFT, because the specific values of the non-zero trainable entries in $\bm{R}$ are irrelevant; only their positions and presence are significant.

A straightforward fully-connected scheme linking two layers requires $\mathcal{O}(d^2)$ edges (corresponding to a single dense orthogonal matrix), specifically $d^2-d$ edges for self-loop exclusion (when $d=4$, this equates to 12 edges). Figure~\ref{fig:info_trans} illustrates a possible matrix factorization, which involves only $10$ edges—fewer than those in a conventional dense orthogonal matrix. This perspective offers a graphical lens for analyzing the parameter-efficiency of OFT, making it possible to design practical factorizations through such a representation. The design draws upon the structure of butterfly graphs~\cite{cooley1965algorithm}, introduced by the Cooley-Tukey algorithm, which can interconnect $d$ source and $d$ target nodes in a dense manner using only $\mathcal{O}(d\log d)$ edges. To illustrate, while the arrangement in Figure~\ref{fig:info_trans} achieves full connectivity with $10$ edges, the butterfly configuration attains the same with just $8$ edges. In what follows, we outline how employing the butterfly structure can enhance parameter efficiency.

\section{Orthogonal Parameterization by Butterfly Factorization}

Initially utilized in the Cooley-Tukey algorithm for efficient computation of the fast Fourier transform, the butterfly structure allows localized changes in the frequency domain to affect the spatial domain globally—a phenomenon that parallels the information transmission challenge under consideration, wherein each node at the first layer is able to propagate information to every node at the final layer. Additionally, butterfly topologies are widely adopted in computer network design~\cite{solihin2015fundamentals,li2011linear} for optimized information transfer. Taking $k\geq 2$ as a power of $2$, the \emph{butterfly factor} $\bm{B}^F(k)$ is defined by

\begin{equation}\label{eq:but_factor}
\begin{aligned}
    \bm{B}^F(k)=\begin{bmatrix}
    \text{diag}(\bm{d}_1) & \text{diag}(\bm{d}_2) \\
    \text{diag}(\bm{d}_3) & \text{diag}(\bm{d}_4)
    \end{bmatrix}\in\mathbb{R}^{k\times k},
\end{aligned}
\end{equation}

where each $\bm{d}_i$ is a vector in $\mathbb{R}^{\frac{k}{2}}$. Now, for $d=2^N$, we construct the $d$-dimensional \emph{butterfly matrix} $\bm{B}(d)\in\mathbb{R}^{d\times d}$ recursively as

\begin{equation}
\begin{aligned}
    \!\!\bm{B}(d)=\tilde{\bm{B}}(d,d)\!\cdot\!\begin{bmatrix}
    \bm{B}_1(\frac{d}{2})\!\!\!\!\! & \bm{0} \\
    \bm{0}\!\!\!\!\! &\bm{B}_2(\frac{d}{2})
    \end{bmatrix}= \tilde{\bm{B}}(d,d)\tilde{\bm{B}}(d,\frac{d}{2})\cdots\tilde{\bm{B}}(d,2),
\end{aligned}
\end{equation}

In this context, $\bm{B}_1(\frac{d}{2})$ and $\bm{B}_2(\frac{d}{2})$ denote two butterfly matrices each with dimension $\frac{d}{2}$. The concept of a \emph{butterfly component} is introduced as $\thickmuskip=2mu \medmuskip=2mu \tilde{\bm{B}}(d,k)=\text{diag}(\bm{B}^F_1(k),\cdots,\bm{B}^F_{d/k}(k))$, which represents a block-diagonal matrix of size $d\times d$ constructed with block size $k$. Here, every $\bm{B}^F_i(k)$ refers to the butterfly factors as specified in Equation~\ref{eq:but_factor}. Utilizing these structures, an orthogonal matrix can be parameterized by ensuring each multiplicative component $\tilde{\bm{B}}(d,k)$ within the butterfly matrix $\bm{B}(d)$ retains orthogonality. Focusing on the block-diagonal matrix $\tilde{\bm{B}}(d,2)$ with block size $2$, orthogonality can be straightforwardly enforced via the Cayley transform (or equivalently, by using $2$-dimensional rotations) to parameterize each constituent block, following approaches in \cite{liu2021orthogonal,qiu2023controlling}. Each butterfly component exhibits a non-zero structure that is a permutation of the non-zero arrangement present in $\tilde{\bm{B}}(d,2)$, facilitating a straightforward orthogonal parameterization of all butterfly components. This method provides an efficient way to represent orthogonal matrices based on numerous $2\times 2$ orthogonal building blocks. Building on \cite{chen2022pixelated}, we extend the idea of butterfly matrices by introducing a block butterfly component $\tilde{\bm{B}}^b(d,k)$, in which every element $\bm{d}_i$ is now replaced by a $b\times b$ matrix. To assure the orthogonality of the block butterfly component $\tilde{\bm{B}}^b(d,2)$, each block of size $2b\times 2b$ is independently parameterized as an orthogonal matrix. For $\tilde{\bm{B}}^b(d,k)$ with $k>2$, their non-zero configurations amount to blockwise permutations of that in $\tilde{\bm{B}}^b(d,2)$, thus enabling similar orthogonal constructions. Synthesizing these elements, the BOFT forward pass is then expressed as

\begin{equation}
    \begin{aligned}\nonumber
    \bm{z}=\big(\bm{R}(m,b)\cdot\bm{W}^0\big)^\top\bm{x},~~~~\text{s.t.}~\bigg{\{}\bm{R}(m,b)=\prod_{i=1}^m \tilde{\bm{B}}^b_{(d,i)}~~\&~~\underbrace{\big(\tilde{\bm{B}}^b_{(d,j)}\big)^\top\tilde{\bm{B}}^b_{(d,j)}=\tilde{\bm{B}}^b_{(d,j)}\big(\tilde{\bm{B}}^b_{(d,j)}\big)^\top\!=\bm{I}_d}_{\forall j\in [1, m]}\bigg{\}},
    \end{aligned}
\end{equation}

Here, we introduce the notation $\tilde{\bm{B}}^b(d,2^{m - i+1})$ as $\tilde{\bm{B}}^b_{(d,i)}$ for convenience. $\bm{I}_d$ denotes the $d\times d$ identity matrix. The orthogonal matrix $\bm{R}(m,b)\in\mathbb{R}^{d\times d}$ is constructed as a product of $m$ orthogonal butterfly blocks. For clarity, we refer to BOFT employing $\bm{R}(m,\frac{b}{2})$ as BOFT($m$,$b$), provided $b\geq 2$. When $m=1$, BOFT($1$,$b$) simplifies to the block-diagonal OFT~\cite{qiu2023controlling} with block size $b$. Further, if $b=d$, BOFT($1$,$d$) matches the original OFT~\cite{qiu2023controlling}, which involves a full unconstrained orthogonal matrix. When $m = \log \frac{2d}{b}$, BOFT utilizes the block butterfly matrix $\bm{B}^{\frac{b}{2}}(d)$ as $\bm{R}$, leading to a fully dense orthogonal $\bm{R}$. In the general case, the parameterization of BOFT($m$,$b$) provides $\frac{1}{2}(b-1)dm$ trainable parameters for adapting a $d\times n$ linear layer. In the special scenario where the butterfly matrix is used, that is, $m=\log d, b=2$, the number of parameters for BOFT scales as $\mathcal{O}(d\log d)$. By contrast, the original OFT with a dense orthogonal matrix requires $\mathcal{O}(d^2)$ parameters, while block-diagonal OFT using $r$ blocks needs $\mathcal{O}(bd)$ parameters. Consequently, to obtain a dense orthogonal matrix, the standard OFT must choose $b=d$ for its blocks, whereas BOFT does not require such restriction on $b$ to produce density.

\textbf{Identity initialization for BOFT}. Fine-tuning routines generally leverage the weights of a pretrained model directly, so that the changes made during finetuning remain close to the starting point. For instance, LoRA initializes its low-rank updates with zeros. Analogously, in BOFT, we set all butterfly factors to identity matrices at initialization (that is, initializing the skew-symmetric matrix to zeros in the Cayley parameterization).

\textbf{Multiplicative Dropout for BOFT}. In LoRA~\cite{hulora2022}, dropout is incorporated for the low-rank updates to combat overfitting. Standard dropout~\cite{srivastava2014dropout} can be readily employed in LoRA, but cannot be directly used in BOFT because of the multiplicative manner in which the weights are updated. To resolve this limitation, we introduce a multiplicative dropout tailored for BOFT. Since $\bm{R}(m,b)$ comprises $m$ butterfly layers, each being transformable into $2b\times 2b$ block-diagonal orthogonal matrices, our multiplicative dropout mechanism randomly selects a proportion $p_1$ of these butterfly layers and, within each selected layer, a proportion $p_2$ of the diagonal blocks, replacing those blocks by identity matrices.

\textbf{Expressivity of BOFT}. The butterfly structure, when combined with appropriate permutations, is capable of exactly representing numerous well-known fast linear transforms~\cite{dao2019learning,dao2019kaleidoscope} (\emph{e.g.}, fast Fourier transform, Hadamard transform). Nevertheless, it is still an open question how closely our orthogonal butterfly matrix can approximate an arbitrary orthogonal matrix. To explore this, we carried out an experiment approximating a randomly sampled dense orthogonal matrix~\cite{anderson1987generation} of dimension $512\times 512$. The outcomes, depicted in Figure~\ref{fig:para_eff}, are averaged across 10 random initializations. The y-axis reflects the approximation error, while the x-axis represents the effective number of tunable parameters. Each curve of the same color corresponds to BOFT configurations with identical block sizes; specifically, the leftmost point corresponds to BOFT($1$,$b$) (\emph{i.e.}, a standard block-diagonal OFT with block size $b$). In general, BOFT demonstrates superior parameter efficiency relative to OFT. For instance, BOFT($9$,$2$) achieves higher expressivity than BOFT($1$,$16$) despite using significantly fewer parameters. Moreover, a configuration with smaller $b$ and larger $m$ tends to achieve greater parameter efficiency; for example, BOFT($6$,$4$) attains a similar approximation error as BOFT($2$,$16$) with many fewer parameters. Overall, the butterfly matrix encodes a richer subset of the orthogonal group than the block-diagonal construction, making BOFT inherently more expressive than OFT at a fixed block size.

\begin{theorem}[Expressivity of BOFT]\label{thm:exp_boft}
    BOFT exhibits strictly greater expressivity than OFT given the same block size. The full range of orthogonal matrices of size $d$ can be approximated by compositions such as $\bm{B}_{d-1,1}(d)\bm{B}^\top_{d-1,2}(d)\cdots\bm{B}_{1,1}(d)\bm{B}^\top_{1,2}(d)$, where all $\bm{B}_{i,j}(d)$ denote butterfly matrices.
\end{theorem}

According to Theorem~\ref{thm:exp_boft}, BOFT admits a straightforward extension: the final orthogonal transformation can be generalized as $\thickmuskip=2mu \medmuskip=2mu \bm{R}^G(m_1,b_1,m_2,b_2,l)=\bm{R}_{l,1}(m_1,b_1)\bm{R}_{l,2}^T(m_2,b_2)\cdots\bm{R}_{1,1}(m_1,b_1)\bm{R}_{1,2}^T(m_2,b_2)$, where $\bm{R}_{i,j}^T(m,b)$ signifies the orthogonal matrices utilized within BOFT. When we select $m_1=m_2=\log d$, $b_1=b_2=2$, and $l=d-1$, the resulting $\bm{R}^G(m_1,b_1,m_2,b_2,l)$ spans the complete orthogonal group; this particular form is referred to as the kaleidoscope hierarchy~\cite{dao2019kaleidoscope}. Despite this higher expressivity, it should be noted that greater expressiveness does not invariably yield better finetuning outcomes—full finetuning, while theoretically fully expressive, can in practice give unsatisfactory results. Thus, achieving a suitable balance between expressivity and regularity is fundamental for robust finetuning and generalization. BOFT enhances the space of tunable parameters by introducing structured priors, which in turn facilitates discovery of an optimal trade-off between flexibility and regularity.

\textbf{Spectral properties}. Compared to LoRA, orthogonal finetuning demonstrates superior preservation of spectral characteristics, as it exactly maintains the spectral norm of the pretrained weight matrix $\bm{W}^0$. This can be illustrated using the singular value decomposition: $\bm{W}^0=\bm{U}\bm{\Sigma}\bm{V}^\top$, where $\bm{U}$ and $\bm{V}$ are orthogonal matrices, and $\bm{\Sigma}$ contains the singular values along its diagonal. In both OFT and BOFT, an orthogonal matrix $\bm{R}$ is applied from the left, resulting in finetuned weights $\bm{R}\bm{U}\bm{\Sigma}\bm{V}^\top$. This left multiplication by an orthogonal matrix leaves the largest singular value unchanged, thereby preserving the spectral norm of $\bm{W}^0$. Prior studies have demonstrated that this property significantly enhances both the stability of training and the model's ability to generalize~\cite{miyato2018spectral,yoshida2017spectral}. Further mathematical details regarding these properties can be found in Appendix~\ref{app:sn_property}.

\textbf{Orthogonal finetuning as learning bilinear similarity}. BOFT may be interpreted as optimizing the bilinear similarity function $\bm{w}^0_i\bm{R}\bm{x}$, where $\bm{w}^0_i$ denotes the $i$-th neuron (specifically, the $i$-th column) from the pretrained weight matrix $\bm{W}^0$. Within this formulation, BOFT effectively learns a bilinear similarity matrix $\bm{R}$, subject to a strong regularization constraint—namely, that $\bm{R}$ must be orthogonal. This links BOFT conceptually to the field of distance metric learning~\cite{xing2002distance} and also aligns with the broader mathematical framework of bilinear forms~\cite{roman2005advanced}.

\textbf{Inductive bias and generalization in BOFT}. The matrix $\bm{R}(m,b)$ employed in BOFT often denotes a structurally restricted subset of the broader orthogonal group, which imposes constraints on the class of admissible models. Consequently, BOFT introduces a particular inductive bias by design. We contend that the specific structured bias imparted by butterfly factorization can promote improved generalization, as it incorporates structure present in several foundational linear transforms~\cite{dao2019kaleidoscope}, including the discrete Fourier, sine/cosine, and Hadamard transforms. Additionally, the inherent sparsity of matrix factorization in BOFT could introduce further implicit inductive biases~\cite{gunasekar2017implicit,li2018algorithmic,liu2019neural}.

\textbf{Comparison to butterfly-based sparse training}. Several prior studies~\cite{dao2019learning,dao2019kaleidoscope,chen2022pixelated,dao2022monarch} have investigated the use of the butterfly parameterization in sparse training contexts. These works predominantly employ the butterfly structure as a direct reparameterization for weight matrices and train neural architectures starting from randomly initialized parameters. In contrast, \cite{dao2022monarch} explores an alternative by adapting pretrained weights: the approach involves projecting these weights onto a butterfly matrix variant, followed by optimization of these projected elements for specific downstream objectives. BOFT, however, introduces a fundamentally distinct approach to finetuning—modifying weights through transformation with weight matrices that are shared across layers.

\input{rebuttal-results}

\section{Concluding Remarks and Limitations}

This work introduces Orthogonal Butterfly, a broadly applicable and parameter-efficient finetuning technique for foundation models, underpinned by the butterfly matrix structure. Our primary conceptual contribution is to enhance parameter efficiency by expressing a dense orthogonal matrix as a product of multiple sparse orthogonal matrices. To facilitate the discovery of valid matrix factorizations, we put forward a framework centered on graphical information transmission. Through this perspective, we demonstrate that the butterfly construction aligns well with our goals for sparse orthogonal matrix decompositions. Empirical results support the practical utility of BOFT across finetuning scenarios in large language models, large-scale vision models, and text-to-image generative models. These experiments highlight the advantage of BOFT as a versatile and robust approach for model finetuning.

Nonetheless, BOFT does have certain drawbacks. Because BOFT's final orthogonal transformation arises from multiplying several orthogonal matrices, it incurs a somewhat higher training time overhead than OFT. The challenge of reducing BOFT's training overhead remains unresolved. Positively, post-finetuning, the orthogonal matrices learned via BOFT can be pre-composed into the pretrained model, thus avoiding any inference-time latency. Additionally, it is not yet clear whether the butterfly network represents the optimal structure for information transmission. Our information transmission framework paves the way to explore ideas from the field of computer networking, where the efficiency of network topologies for data transfer is extensively researched. We anticipate that alternative, potentially more efficient, network structures could inform future methods of assembling dense orthogonal matrices.